% Nouveau EGA V (cote 25): archive cote 25, site page 157 (PDF page 158).
% typescript. Draft compared: none.
% AI-transcribed by Claude Opus 5.5 (claude-opus-5-5, adaptive thinking, effort high), Message Batches API, batch ega25b (msgbatch_01A7n5YuXubUPNK1zsepi8RG), 2026-11-13; not yet checked by eye.
\scanpage{25/157}
% NOTE: A complete forward-reading pale impression is visible. Its physical origin (carbon or material visible through a sheet) cannot be settled from these images. The recoverable text is transcribed rather than omitted; the exceptionally faint readings remain doubtful. The neighbouring source transcription was consulted only as a comparison proposal, not as evidence that a word is legible here.
\var{\margin{\unsure{nouveau EGA V}}}{\placedmargin{\unsure{nouveau EGA V}}}{marginal note placed}{ega25_25_157-v0001}
\unsure{IV 17.15.} \uline{\unsure{Formes lisses}}; \var{\inserted{\unsure{augmentations et singularités élémentaires.}}}{\unsure{augmentations et singularités élémentaires.}}{insertion}{ega25_25_157-v0002} \var{\deleted{\var{\xout{\unsure{x}}}{}{typewriter strike-out}{ega25_25_157-v0004} \unsure{formes quadratiques lisses}}}{}{deletion}{ega25_25_157-v0003}
% NOTE: Pale curved traces at the upper left follow the sheet edge; their material or writing origin is unresolved, so no textual bracket is assigned.
\unsure{NB Je viens de m'apercevoir que la terminologie introduite dans}
\unsure{ma rédaction} \var{\inserted[typed]{\unsure{SUP}}}{\unsure{SUP}}{insertion}{ega25_25_157-v0005} \unsure{de la Supersingularité est déraisonnable, et conflicte}
\unsure{en particulier avec la terminologie reçue. (De toutes façons, tu a}\var{\xout{\unsure{s}}}{}{typewriter strike-out}{ega25_25_157-v0006} \unsure{dû}
\unsure{no}\var{\over{\unsure{t}}{\unsure{t}}}{\unsure{t}}{writing over}{ega25_25_157-v0007}\unsure{er que dans la déf.1 de ce texte, il faut lire ``dégénérée'' au lieu de}
\unsure{``non dégénérée''). Le canular provient des variétés de dimension paire}
\unsure{en car.2,} \var{\xout{\unsure{une telle variété}}}{}{typewriter strike-out}{ega25_25_157-v0008} \unsure{définie}\hand{\unsure{s}} \unsure{par exemple par une équation} $\hand{\unsure{\phi}}\unsure{=0}$
\unsure{dans une variété ambiante de dim impaire, car avec la terminologie de mes}
\unsure{notes, une telle variété ne peut avoir de} \var{\xout{\unsure{zéro sing}}}{}{typewriter strike-out}{ega25_25_157-v0009} \hand{\unsure{``}}\unsure{point singulier}
\unsure{or}\var{\over{\unsure{d}}{\unsure{d}}}{\unsure{d}}{writing over}{ega25_25_157-v0010}\unsure{inaire}\hand{\unsure{''}} \unsure{(= singularité quadratique ordinaire), i.e.} $\hand{\unsure{\phi}}$ \unsure{ne peut avoir}
\unsure{de} \hand{\unsure{``}}\unsure{zéro singulier ordinaire}\hand{\unsure{''}}, \unsure{dû au fait qu'en car 2, une forme quadratique}
\unsure{à un nombre impair de variables est toujours ``dégénérée''. Or tout le}
\unsure{monde a toujours considéré que même en car. 2, l'origine du cône affine}
$\unsure{x^2+yz=0}$ \unsure{est une singularité quadratique ordinaire. Dans les présen-}%
\unsure{tes notes, je dégage la notion de} \var{\xout{\unsure{f}}}{}{typewriter strike-out}{ega25_25_157-v0011}\unsure{forme quadratique lisse (ou ``ordinaire''),}
\unsure{(sur un Module loc libre de type fini), de telle façon que}
\unsure{non dégénéré} \hand{$\unsure{\Longrightarrow}$} \unsure{lisse, la réciproque étant vraie si E est de rang pair}
\uline{\unsure{ou}} \unsure{les car résiduelles de la base S sont toutes} $\hand{\unsure{\neq}}\unsure{2}$. \unsure{Cette notion étant}
\unsure{posée, je propose de garder la terminologie ``supersingulière'' (qui ne}
\unsure{conflicte avec aucune terminologie reçue)} \var{\xout{\illegible{2} \unsure{de mon texte}}}{}{typewriter strike-out}{ega25_25_157-v0012} \unsure{de SUP,}
\unsure{comme correspondant à la notion de forme quadratique non dégéné}\var{\over{\unsure{f}}{\unsure{r}}}{\unsure{r}}{writing over}{ega25_25_157-v0013}\unsure{ée; on}
\unsure{parlera de même de point supersingulier d'un préschéma loc noeth., ou}
\unsure{géométriquement supersingulier pour un préschéma loc de type fini sur un}
\unsure{corps, dans le même sens (en regardant l'hom} $\unsure{\mathrm{Sym}(\underline{M}/\underline{m}^2)}\hand{\unsure{\longrightarrow}}\unsure{\mathrm{gr}_{\underline{m}}(A)}$ \unsure{et}
\unsure{en disant que le noyau ne peut être engendré par une forme quadratique}
\unsure{non dégénérée); donc ce qui dans SUP s'appelait à tort ``zéro singulier}
\unsure{ordinair}\hand{\unsure{e}}\unsure{'' s'appelera ``zéro singulier non supersingulier'', ou ``zéro}
\mspage{\hand{157}}
